\subsection*{Câu 1}

\noindent\textbf{1. Mạch kiến thức:} Thứ tự thực hiện phép tính

\noindent\textbf{2. Độ khó:} **

\noindent\textbf{3. Đề bài:}
\[
12+3\cdot 2^5:4-3
\]

\noindent\textbf{4. Lời giải:}
\begin{align*}
12+3\cdot 2^5:4-3
    &= 12+6^5:4-3 \\
    &= 12+7776:4-3 \\
    &= 12+1944-3 \\
    &= 1953
\end{align*}

\noindent\textbf{5. Token bắt đầu sai:}
\[
6
\]

\noindent\textbf{6. Giải thích lỗi:}

Theo thứ tự thực hiện phép tính, lũy thừa được thực hiện trước phép nhân và phép chia.

\noindent\textbf{7. Sửa lại:}

\textit{Dạng câu hỏi:} Trắc nghiệm

\textit{Câu hỏi:} Phép biến đổi nào sau đây được thực hiện đúng từ biểu thức
\[
12+3\cdot 2^5:4-3=...
\]

\begin{itemize}
    \item[A.] 
    \[
    15\cdot 2^5:4-3
    \]
    
    \item[B.] 
    \[
    12+3\cdot 32:1
    \]
    
    \item[C.] 
    \[
    12+3\cdot 32:4-3
    \]
    
    \item[D.] 
    \[
    12+3\cdot 10:4-3
    \]
\end{itemize}

\textit{Đáp án:} C

\subsection*{Câu 2}

\noindent\textbf{1. Mạch kiến thức:} Quy tắc dấu ngoặc

\noindent\textbf{2. Độ khó:} **

\noindent\textbf{3. Đề bài:}
\[
92-(55-8)+(-45)
\]

\noindent\textbf{4. Lời giải:}
\begin{align*}
92-(55-8)+(-45)
    &= 92-55+8+45 \\
    &= (92+8)+(-55+45) \\
    &= 90+10 \\
    &= 100
\end{align*}

\noindent\textbf{5. Token bắt đầu sai:}
\[
+45
\]

\noindent\textbf{6. Giải thích lỗi:}

Khi bỏ dấu ngoặc có dấu ``$+$'' đằng trước, ta giữ nguyên dấu của các số hạng trong ngoặc.

\noindent\textbf{7. Sửa lại:}

\textit{Dạng câu hỏi:} Trắc nghiệm

\textit{Câu hỏi:} Tính một cách hợp lí, kết quả biểu thức là:

\begin{itemize}
    \item[A.] $-10$
    \item[B.] $0$
    \item[C.] $10$
    \item[D.] $100$
\end{itemize}

\textit{Đáp án:} B

\subsection*{Câu 3}

\noindent\textbf{1. Mạch kiến thức:} Thứ tự thực hiện phép tính với số hữu tỉ

\noindent\textbf{2. Độ khó:} ***

\noindent\textbf{3. Đề bài:}
\[
\left\{
-\frac{2}{3}
-
\left[
-\frac{1}{18}:\left(\frac{1}{3}\right)^2
\right]
\right\}\cdot\frac{5}{6}
\]

\noindent\textbf{4. Lời giải:}
\begin{align*}
\left\{
-\frac{2}{3}
-
\left[
-\frac{1}{18}:\left(\frac{1}{3}\right)^2
\right]
\right\}\cdot\frac{5}{6}
&=
\left\{
-\frac{2}{3}
-
\left[
\left(-\frac{1}{6}\right)^2
\right]
\right\}\cdot\frac{5}{6} \\
&=
\left\{
-\frac{2}{3}
-
\left[
\frac{1}{36}
\right]
\right\}\cdot\frac{5}{6} \\
&=
\left(
-\frac{2}{3}-\frac{1}{36}
\right)\cdot\frac{5}{6} \\
&=
-\frac{25}{36}\cdot\frac{5}{6} \\
&=
-\frac{125}{216}
\end{align*}

\noindent\textbf{5. Token bắt đầu sai:}
\[
-\frac{1}{6}
\]


\noindent\textbf{6. Giải thích lỗi:}

Theo thứ tự thực hiện phép tính, lũy thừa được thực hiện trước phép nhân và phép chia.

\noindent\textbf{7. Sửa lại:}

\textit{Dạng câu hỏi:} Trắc nghiệm

\textit{Câu hỏi:} Chọn phép biến đổi đúng:
\[
\left\{
-\frac{2}{3}
-
\left[
-\frac{1}{18}:\left(\frac{1}{3}\right)^2
\right]
\right\}\cdot\frac{5}{6}=....
\]

\begin{itemize}
    \item[A.]
    \[\left\{-\frac{2}{3}-\left[-\frac{1}{18}:\frac{1}{9}\right] \right\}\cdot\frac{5}{6}\]

    \item[B.]\[\left\{ -\frac{2}{3} - \left[\frac{1}{18}:\frac{1}{3}\right]^2 \right\}\cdot\frac{5}{6}\]

    \item[C.]\[ \left\{-\frac{2}{3}-\left[-\left(\frac{1}{6}\right)^2\right]\right\}\cdot\frac{5}{6}  \]

    \item[D.]
    \[\left\{ -\frac{2}{3} -\left[-\frac{1}{18}:\frac{1}{3}\right]\right\}\cdot\frac{5}{6}\]
\end{itemize}

\textit{Đáp án:} A

\subsection*{Câu 4}

\noindent\textbf{1. Mạch kiến thức:} Thứ tự thực hiện các phép tính. Quy tắc chuyển vế

\noindent\textbf{2. Độ khó:} **

\noindent\textbf{3. Đề bài:}
\[
1-2x=\frac{7}{5}-\frac{2}{5}:\frac{2}{3}
\]

\noindent\textbf{4. Lời giải:}
\begin{align*}
1-2x
    &= \frac{7}{5}-\frac{2}{5}:\frac{2}{3} \\
1-2x
    &= \frac{5}{5}:\frac{2}{3} \\
1-2x
    &= 1:\frac{2}{3} \\
1-2x
    &= \frac{3}{2} \\
-2x
    &= \frac{3}{2}-1 \\
-2x
    &= \frac{1}{2} \\
x
    &= -\frac{1}{4}
\end{align*}

\noindent\textbf{5. Token bắt đầu sai:}
\[
\frac{5}{5}
\]

\noindent\textbf{6. Giải thích lỗi:}

Trong biểu thức có phép cộng, trừ, nhân, chia, phải thực hiện phép nhân, chia trước phép cộng, trừ.

\noindent\textbf{7. Sửa lại:}

\textit{Dạng câu hỏi:} Trắc nghiệm

\textit{Câu hỏi:} Chọn biểu thức đúng để thay thế cho vế phải.

\begin{itemize}
    \item[A.] $1-2x=1:\frac{2}{3}$
    \item[B.] $1-2x=\frac{7}{5}-\frac{4}{15}$
    \item[C.] $1-2x=\frac{7}{5}-\frac{2}{5}$
    \item[D.] $1-2x=\frac{7}{5}-\frac{3}{5}$
\end{itemize}

\textit{Đáp án:} D

\subsection*{Câu 5}

\noindent\textbf{1. Mạch kiến thức:} Quy tắc chuyển vế

\noindent\textbf{2. Độ khó:} ***

\noindent\textbf{3. Đề bài:}
\[
\left(-\frac{5}{6}x+\frac{5}{4}\right):\frac{3}{2}
=\frac{4}{3}
\]

\noindent\textbf{4. Lời giải:}
\begin{align*}
\left(-\frac{5}{6}x+\frac{5}{4}\right):\frac{3}{2}
    &= \frac{4}{3} \\
-\frac{5}{6}x+\frac{5}{4}
    &= \frac{4}{3}\cdot\frac{3}{2} \\
-\frac{5}{6}x+\frac{5}{4}
    &= 2 \\
-\frac{5}{6}x
    &= 2+\frac{5}{4} \\
-\frac{5}{6}x
    &= \frac{13}{4} \\
x
    &= -\frac{39}{10}
\end{align*}

\noindent\textbf{5. Token bắt đầu sai:}
\[
+
\]

\noindent\textbf{Dòng bắt đầu sai:} 4

\noindent\textbf{6. Giải thích lỗi:}

Trong một đẳng thức, khi chuyển một số hạng từ vế này sang vế kia, phải đổi dấu của số hạng đó.

\noindent\textbf{7. Sửa lại:}

\textit{Dạng câu hỏi:} Trắc nghiệm

\textit{Câu hỏi:} Chọn phép biến đổi đúng để tiếp tục tìm $x$
\[
-\frac{5}{6}x+\frac{5}{4}
    = 2 
\]

\begin{itemize}
    \item[A.] $-\frac{5}{6}x=2+\frac{5}{4}$
    \item[B.] $-\frac{5}{6}x=2-\frac{5}{4}$
    \item[C.] $-\frac{5}{6}x=-2-\frac{5}{4}$
    \item[D.] $-\frac{5}{6}x=-2+\frac{5}{4}$
\end{itemize}

\textit{Đáp án:} B